\section{SCL Development and Refinement}\label{sec:development}

This section presents the development and refinement of \scl{} using the development data
described in Sect.~\ref{sec:data}. Section~\ref{sec:prior} introduces the initial always-on banding
(\aob{}) design and discusses its limitations. Section~\ref{sec:steps} describes the incremental
improvements that addressed these limitations, supported by the experimental comparisons in
Table~\ref{tab:ablation}. Finally, Sect.~\ref{sec:freeze} explains the configuration freezing and
validation procedures performed before external evaluation to ensure reproducibility and prevent
changes based on subsequent results.

\begin{table*}[!tp]
\caption{Design steps from \aob{} to \scl{} on development
data (HOTA, \%), and versions that were tested and rejected. The floor of
ByteTrack and OC-SORT is given in parentheses.}
\label{tab:ablation}
\centering\footnotesize
\setlength{\tabcolsep}{4pt}
\begin{tabular}{@{}>{\raggedright\arraybackslash}p{25mm}>{\raggedright\arraybackslash}p{36mm}>{\raggedright\arraybackslash}p{62mm}@{}}
\toprule
Step & Change & Effect\\
\midrule
\aob{} emulated & always-noisy bands, overlap deduplication, rank remap & ByteTrack (0.1) 56.41 vs host 67.70; OC-SORT (0.1) 45.72 vs 66.44\\
+ mode gate & mode statistic, host-anchored clean mode, crowd-safe duplicates & SparseTrack 64.72 to 68.93 (host 68.88); BoostTrack 62.61 to 68.37 (host 68.49)\\
+ whole-stream mode & mode from all past frames & ByteTrack (0.1) 67.561 to 67.702; OC-SORT (0.1) 65.697 to 66.071\\
+ continuation rescue & foreground track-consistent rescue & OC-SORT (0.01) 66.405 to 66.619; OC-SORT (0.1) 66.071 to 66.469\\
+ background check = \scl{} & no background mode gives no evidence against the host & ByteTrack/BoostTrack (0.1) identical to host; OC-SORT 67.041 (0.01), 66.907 (0.1)\\
\midrule
\multicolumn{2}{@{}l}{\textit{Rejected version}} & \textit{Evidence}\\
\multicolumn{2}{@{}l}{no candidates at frame 1} & $-0.3$ to $-0.5$ HOTA on every MOT17 host\\
\multicolumn{2}{@{}l}{rescue below a two-stage host's low stage} & FP $+700$ to $+1000$, $-0.4$ to $-0.5$ HOTA\\
\multicolumn{2}{@{}l}{rank-remapped scores in clean frames} & ByteTrack $-0.2$ / $-1.6$ HOTA\\
\multicolumn{2}{@{}l}{host-clipped noisy primary band} & KITTI YOLOv8n ByteTrack $+1.4$ HOTA but uncontrolled false tracks on a noisy aerial stream (26.9 vs 14.9 boxes/frame)\\
\bottomrule
\end{tabular}
\end{table*}

\begin{table}[!tb]
\caption{Baseline comparison of the development stage on MOT17 val-half
(published YOLOX-X detections). Host: the tracker alone, without any
layer (our run of the published tracker). $\Delta$: tracker with the earlier design \aob{} or with \scl{}
minus host, with 95\% paired-bootstrap interval (10,000 resamples,
seed 42). Bold intervals exclude zero. IDS: identity switches, a count
(lower is better).}
\label{tab:prior}
\centering\footnotesize
\setlength{\tabcolsep}{3pt}
\renewcommand{\arraystretch}{1.1}
\begin{tabular}{@{}lrcc@{}}
\toprule
Metric & Host & $\Delta$ \aob{} & $\Delta$ \scl{}\\
\midrule
\multicolumn{4}{@{}l}{\textit{SparseTrack} \cite{liu2025sparsetrack}}\\
\quad HOTA & 68.88 & \cibt{-4.15}{-5.54}{-1.66} & \cit{+0.06}{-0.01}{+0.25}\\
\quad MOTA & 77.85 & \cibt{-6.14}{-8.10}{-2.74} & \cit{+0.08}{-0.03}{+0.30}\\
\quad IDF1 & 81.97 & \cibt{-4.49}{-6.19}{-1.96} & \cit{+0.16}{-0.02}{+0.70}\\
\quad IDS & 124 & \cit{-16}{-47}{+12} & \cit{+7}{0}{+19}\\
\multicolumn{4}{@{}l}{\textit{BoostTrack, online} \cite{stanojevic2024boosttrack}}\\
\quad HOTA & 68.49 & \cibt{-5.89}{-7.76}{-2.87} & 0 (identical)\\
\quad MOTA & 75.50 & \cibt{-8.87}{-12.37}{-7.14} & 0 (identical)\\
\quad IDF1 & 81.41 & \cibt{-6.22}{-9.33}{-2.85} & 0 (identical)\\
\quad IDS & 113 & \cibt{-38}{-81}{-1} & 0 (identical)\\
\multicolumn{4}{@{}l}{\textit{BoostTrack, with GBI} \cite{stanojevic2024boosttrack}}\\
\quad HOTA & 71.72 & \cibt{-5.37}{-7.30}{-2.75} & 0 (identical)\\
\quad MOTA & 81.03 & \cibt{-8.73}{-12.59}{-6.50} & 0 (identical)\\
\quad IDF1 & 84.16 & \cibt{-5.65}{-8.96}{-2.77} & 0 (identical)\\
\quad IDS & 97 & \cit{-27}{-73}{+11} & 0 (identical)\\
\bottomrule
\end{tabular}
\end{table}

\subsection{Initial Design and Limitations}\label{sec:prior}
The earlier design \aob{} computed the same nested Otsu groups but treated every stream as noisy:
in every frame it applied the groups, removed every weaker candidate that overlapped a stronger one
with IoU above one half, and rescored all candidates by rank. Table~\ref{tab:prior} compares \aob{}
and \scl{} with strong published trackers on MOT17. \aob{} lowered the HOTA of SparseTrack by 4.15
points and of BoostTrack by 5.89 points online and 5.37 points with its gap-filling post-processing
(GBI), for two reasons: removing overlapping boxes deleted people hidden behind others in crowds,
and its thresholds were stricter than the trackers' own. With the same trackers, \scl{} changed
HOTA by $+0.06$ for SparseTrack and left both BoostTrack versions identical to the host. A layer
that acts in every frame therefore cannot be placed safely in front of an unknown tracker, and the
next steps let the layer decide from the scores whether to act at all.


\subsection{Incremental Design Improvements}\label{sec:steps}
Each step was compared with the step before it (Table~\ref{tab:ablation}). The \emph{mode gate}
added the mode check, the clean mode that keeps the tracker's scores and thresholds, and the crowd-safe duplicate
rule; it restored SparseTrack from 64.72 to 68.93 HOTA and BoostTrack from 62.61 to 68.37. The
\emph{whole-stream mode} removed mode changes in the middle of streams and raised ByteTrack from
67.561 to 67.702 and OC-SORT from 65.697 to 66.071. The \emph{continuation rescue} added 0.2 to 0.4
points for OC-SORT. The \emph{background check} completed \scl{}: ByteTrack and BoostTrack returned
to output identical to the host, and OC-SORT rose to 67.04 and 66.91 at the two floors. Four other
variants were tested and not kept.


\subsection{Configuration Freezing and Validation}\label{sec:freeze}
The final configuration was frozen before any external test as a version-control tag, protected
by a lock file with a cryptographic hash of the layer, its configuration, and the evaluation
scripts. Sixty-one automated tests check, among others, that the layer uses no future data and
passes clean streams unchanged. All results of Sect.~\ref{sec:results} come from this frozen
configuration; the comparisons, the ablation, and the sensitivity analysis were run after the
freeze, through run-time options of the evaluation scripts, and changed none of its values.
